\chapter{Implementation}
\label{ch:impl}

This chapter describes how the design is realized. It covers the programming
language and platform, the code conventions, the difficulties met during
development, the dataset and the implementation of each validation layer, and
the algorithms behind the mutators and the trained detector. Core procedures are
shown as code listings, and the full set of mutator transformations is collected
in Appendix~\ref{apndx:code}. The chapter closes with a summary.

%----------------------------------------------------------

\section{Programming Language Selection}

The system is implemented in Python 3.11. The choice follows the domain. The
benchmark analyzes Python code, so the sink patterns, the AST transformations,
and the semantics-preservation checks operate in the same language they target.
The standard \texttt{ast} module provides a stable syntax tree that the mutators
and the negative builder rewrite node by node, and \texttt{ast.unparse}
regenerates source from a modified tree, which makes a round trip from source to
tree and back the basic unit of every transformation. The machine learning stack
the trained detector needs, PyTorch and the HuggingFace Transformers library, is
native to Python, so the whole pipeline lives in one language and one runtime.
Regular expressions back the sink-presence filter, and the \texttt{difflib}-style
line comparison backs the diff filter, both from the standard library.

\section{Platform Selection}

The pipeline and harness run on a single Linux or macOS workstation. The SAST
detectors (Bandit 1.9.4, Semgrep 1.163.0) and the whole-project detectors (Snyk
Code 1.1305.0, CodeQL) run as local subprocesses, and each runner parses the
tool's native output format into a common per-CWE verdict. The LLM detectors
(Claude Sonnet 4.6, GPT-4o, DeepSeek R1) run through their hosted APIs at
temperature zero. The trained detector is fine-tuned on a single CUDA-capable
GPU. Source retrieval uses Git for fix commits and SQLite for the CVEfixes
database.

\section{Code Conventions}

\subsection{Naming Convention}

Modules are named for their role in the pipeline: the scraper, the sink-presence
filter, the audit driver, the diff filter, the negative builder, the splitter,
the mutators, the metrics module, and the detector runners. Each sample carries a
stable identifier that threads through every stage, so a prediction in the
evaluation logs can be traced back to the source file and the validation decision
that admitted it. CWE classes are referred to by their numeric code throughout,
and mutators by their family and index (M1 to M7).

\subsection{Artifact Schema}

Each sample is stored as a paired source file and a schema-versioned metadata
record. The record carries the CWE label, the matched sink pattern, the label
provenance, the CVSS base score and severity, and a label-confidence tier. The
tier separates HIGH-provable advisory positives whose sink line the diff filter
confirmed changed, HIGH-curated canonical samples that are correct by
construction, and MEDIUM-curator \texttt{vudenc} samples that are human-labeled
but carry no paired fix to diff against. Versioning the schema lets the pipeline
evolve the record format without invalidating earlier artifacts, and it makes
every validation decision reproducible from the released files.

\section{Difficulties Encountered and Strategies Used to Tackle Them}

The central difficulty was label noise that the sink filter could not catch: a
sample where the sink is present but is not the vulnerability, because the CVE
fix touched a co-changed file. The strategy was the three-round audit, where each
round measured the false-positive rate on disjoint samples and its findings
motivated the next filter. Round~1 (77 samples) measured 52\% and surfaced six
recurring failure patterns:

\begin{itemize}
  \item \textbf{Sink too broad.} A pattern matched a safe overload of the sink,
    for example a parameterized query counted as SQL injection.
  \item \textbf{Co-changed-file noise.} The sink was present but the fix that
    earned the CWE label modified a different file.
  \item \textbf{Documentation match.} The sink token appeared only in a comment
    or docstring.
  \item \textbf{Pattern conflation.} One regex matched two distinct sinks, one
    safe and one unsafe.
  \item \textbf{Inverted operation.} The matched line was the remediation, not
    the vulnerability.
  \item \textbf{Wrapper abstraction.} The real sink sat behind a project helper
    the pattern could not see.
\end{itemize}

Each pattern motivated a specific tightening, and two CWE classes with systematic
issues were dropped: CWE-434 at 100\% false positive, and CWE-798 with a
single-source bias. Round~2 held at 51\% because new samples surfaced
co-changed-file noise on the untouched classes, which motivated the diff filter.
A second difficulty was proving each mutator safe. The strategy was to skip any
node a transformation could not prove safe, for example a rename that might
collide with an imported name, and to require every mutated sample to re-parse.

\section{Data Set}

Positives are aggregated from five primary sources with differing retrieval
methods and label quality, listed in Table~\ref{tab:sources}. GHSA Advisory DB
dominates because it indexes per-CWE advisories with explicit fix commits. The
\texttt{vudenc} corpus is human-labeled but ships no paired fixed version, so it
cannot be diff-validated and is retained as a lower-confidence tier. Twenty
hand-curated canonical samples are textbook OWASP and SANS exemplars that serve
as calibration points: a detector that misses them is broken regardless of its
CVE-corpus performance.

\begin{table}[htb]
\caption{Positive data sources and counts after label validation. Hard negatives
are constructed safe twins.}
\label{tab:sources}
\centering
\begin{tabular}{lr l}
\toprule
\textbf{Source} & \textbf{Count} & \textbf{Retrieval} \\
\midrule
GHSA Advisory DB    & 384 & per-advisory fix commit \\
\texttt{vudenc}     & 209 & pre-labeled corpus \\
OSV / GHSA GraphQL  &  54 & commit references \\
CVEfixes            &  58 & local SQLite query \\
NVD-targeted        &  33 & NVD API $+$ keyword filter \\
Canonical (curated) &  20 & hand-written exemplars \\
\midrule
Hard negatives      & 429 & constructed safe twins \\
\bottomrule
\end{tabular}
\end{table}

\subsection{Layer 1: Sink-Presence Filter}

For each CWE the filter defines a small set of high-recall regex sink patterns,
and it admits a sample only if at least one pattern matches the source after
comments and docstrings are stripped. The patterns match both safe and unsafe
forms of the sink, so discrimination is left to Layer~3. Two CWEs use proximity
gates: CWE-22 requires an HTTP-request reference within twenty lines of the
sink, because test fixtures with hardcoded paths are otherwise common false
positives. Listing~\ref{lst:sinkfilter} shows the filter in simplified form.

\begin{figure}[H]
\begin{verbatim}
def has_cwe_sink(code, cwe, file_path=None):
    if cwe in TEST_EXCLUSION_CWES and is_test_file(file_path):
        return False, None
    match_code = strip_comments_and_docstrings(code)
    for pat in SINK_PATTERNS[cwe]:
        for m in pat.finditer(match_code):
            if cwe in REQUIRES_REQUEST_PROXIMITY:
                if not has_request_near(code, m.start()):
                    continue
            if cwe in REQUIRES_SECURITY_CONTEXT:
                if not has_security_context_near(code, m.start()):
                    continue
            return True, pat.pattern
    return False, None
\end{verbatim}
\caption{Layer-1 sink-presence filter (simplified). A sample is admitted only if
a CWE sink pattern matches after comment and docstring stripping, subject to the
per-CWE proximity gates.}
\label{lst:sinkfilter}
\end{figure}

Applied to the 1{,}982 pre-filter samples, Layer~1 dropped 749 (38\%). The
rejection rate varied by source from 28\% (NVD) to 68\% (CVEfixes), where
commit-level mislabels were confirmed at scale.

\subsection{Layer 2: Adversarial Audit}

The audit ran three rounds on CWE-stratified samples, each sampling ten items
per populous class and auditing rare classes in full. For each item the auditor
received the identifier, source, repository, file path, the matched regex, and
the code excerpt around the sink, and returned a binary verdict with a
one-sentence justification, without seeing other auditors' verdicts. The audit
is the mechanism that drove the pattern tightenings in the previous section.

\subsection{Layer 3: Diff-Based Filter}

The diff filter implements the evidence-anchored criterion: a positive is valid
only if at least one line containing a sink-pattern match differs between
\texttt{code\_before} and \texttt{code\_after}, after comment stripping and
whitespace normalization. Listing~\ref{lst:difffilter} shows the procedure.

\begin{figure}[H]
\begin{verbatim}
def sink_was_modified(code_before, code_after, cwe):
    patterns = SINK_PATTERNS.get(cwe)
    if not patterns or not code_after.strip():
        return None  # cannot determine

    def sink_lines(code):
        out = set()
        for ln in strip_comments(code).splitlines():
            normalized = " ".join(ln.split())
            if not normalized:
                continue
            for p in patterns:
                if p.search(ln):
                    out.add(normalized)
                    break
        return out

    before = sink_lines(code_before)
    after  = sink_lines(code_after)
    return (before - after) != set()
\end{verbatim}
\caption{Layer-3 diff filter. A positive passes if at least one sink line in
\texttt{code\_before} is absent from \texttt{code\_after}. Whitespace is
normalized so that reformat-only edits do not count as real changes.}
\label{lst:difffilter}
\end{figure}

Of the 738 positives reaching Layer~3, 316 (43\%) had unchanged sink lines and
were rejected; the cut fell unevenly across classes (CWE-79 lost 68\%, CWE-22
lost 71\%). Positives without a paired fixed version, the \texttt{vudenc} tier,
cannot be diff-checked and are retained at the MEDIUM confidence tier, so the
rejected count does not subtract directly to the final positive total. After all
three layers the corpus holds 440 positives. The combined pipeline lowered the
audited false-positive rate from 52\% to 26\%.

\subsection{Final Composition}

The final benchmark is 440 positives plus 429 safe negatives, 869 samples in
all, split 610/127/132 with a repo-group constraint that yields zero cross-split
repository leakage. Table~\ref{tab:composition} gives the per-CWE composition.
Three oversized hard negatives are kept on disk but excluded from active
evaluation, leaving a 129-sample test set (67 positives, 62 negatives).

\begin{table}[htb]
\caption{Final per-CWE composition of the benchmark.}
\label{tab:composition}
\centering
\begin{tabular}{lrrrr}
\toprule
\textbf{Label} & \textbf{Total} & \textbf{Train} & \textbf{Val} & \textbf{Test} \\
\midrule
CWE-89 (SQL Injection)             & 213 & 149 & 32 & 32 \\
CWE-502 (Insecure Deserial.)       &  75 &  52 & 11 & 12 \\
CWE-79 (XSS)                       &  54 &  40 &  7 &  7 \\
CWE-94 (Code Injection)            &  35 &  24 &  5 &  6 \\
CWE-78 (OS Command Injection)      &  26 &  18 &  4 &  4 \\
CWE-918 (SSRF)                     &  22 &  17 &  2 &  3 \\
CWE-22 (Path Traversal)            &  15 &  10 &  2 &  3 \\
\texttt{safe} (hard negatives)     & 429 & 300 & 64 & 65 \\
\midrule
\textbf{Total}                     & \textbf{869} & \textbf{610} & \textbf{127} & \textbf{132} \\
\bottomrule
\end{tabular}
\end{table}

\subsection{Hard-Negative Construction}

The 429 safe samples are not scraped. For each negative the builder takes a real
vulnerable function and applies the canonical remediation for its CWE at the AST
level, for example rewriting \texttt{yaml.load} to \texttt{yaml.safe\_load} or a
string-formatted SQL query to a parameterized query. The remediation preserves
the surface structure and the sink name while neutralizing only the sink
behavior. Each negative therefore carries a provable safe label and is a
minimal-pair counterfactual: a detector that fires on the mere presence of a
sink token is forced into a false positive on the twin.

\section{Algorithms Used}

\subsection{Mutator Transformations}

Each mutator is an AST transformer that walks the parsed tree, rewrites the nodes
it can prove safe, and unparses the result. The transformations are deterministic
given a sample identifier and seed, so the evaluation reproduces. Every mutator
follows three principles: it is semantics-preserving and skips any node it cannot
prove safe; it targets one shortcut, so a per-mutator accuracy drop is
diagnostic; and it is realistic, so identifiers become synonyms and injected code
reads like ordinary scaffolding. The three families attack distinct shortcuts:

\begin{itemize}
  \item \textbf{Sink-preserving (M1 to M4).} Dead-code injection breaks
    positional reasoning, string splitting breaks substring matching, identifier
    rename breaks name memorization, and wrapper extraction breaks
    single-function analysis. All four leave the sink identifier and its
    arguments intact.
  \item \textbf{Sink-targeted (M5, M6).} M5 rewrites an attribute sink
    \texttt{os.system(x)} through a \texttt{getattr} lookup, and M6 rewrites a
    builtin sink \texttt{eval(x)} through a \texttt{builtins} dictionary lookup,
    so in both cases the sink token survives only inside a string literal
    (Appendix~\ref{apndx:code} gives the full rewrites). Both break token-level
    sink matching.
  \item \textbf{Dataflow (M7).} Taint-through-dict routes the sink's first
    argument through a one-entry dict, \texttt{f(\{"a":x\}["a"])}, breaking the
    static lineage from parameter to sink argument while leaving the sink and its
    value unchanged.
\end{itemize}

Three of the seven mutators are applied as online training augmentation, and the
four marked as held-out probes (M2, M5, M6, M7) are used only at test time.
Appendix~\ref{apndx:code} gives a before-and-after listing for each of the seven
mutators.

\subsection{Trained Detector}

The 610-sample training split fine-tunes a dual-task GraphCodeBERT that predicts
both the CWE class and the CVSS base score. Training uses a three-phase schedule.
The first phase is a supervised-contrastive warmup that pulls together
representations of samples sharing a CWE. The second phase applies an LDAM margin
loss \cite{cao2019ldam} on a class-balanced sampler, which widens the decision
margin for the rare classes. The third phase is a deferred re-weighting stage
that switches to Effective-Number class weights \cite{cui2019cb} once the
representation has stabilized. The model
has 126.6M parameters and is treated as a small-data existence proof rather than
a state-of-the-art claim.

\section{Summary}

The implementation realizes the three subsystems in Python: the corpus pipeline
with its three validation layers, shown here as the Layer-1 and Layer-3 listings,
and the constructed negatives; the AST-based mutator suite; and the evaluation
harness with its detector runners and trained model. The dataset tables fix the
composition the next chapters test and evaluate. Chapter~\ref{ch:testing}
describes how the implementation is tested.
